\Needspace{16\baselineskip}
\section{Integrable densities in the compact Gauduchon problem}
\label{sec:integrable-gauduchon-currents}

We keep the Hermitian background $\omega$ and replace the continuous logarithmic density by an integrable density. For an unbounded density, weak compactness of the test currents alone does not pass the determinant integral to the limit; uniform control of the weighted tails provides the missing upper semicontinuity.

We use the following notation from \cref{sec:individual-density-subsolutions,sec:integrable-gauduchon-tests}, with $\omega$ in place of $h$. For a nonzero nonnegative $f\in L^1(\mathrm{d}V)$, the subsolution threshold $b_\omega(f)\geq0$ is, by \cref{thm:individual-density-subsolution}, the largest $b\geq0$ for which some $u\in\PSH(X,\omega)$ satisfies $\mathrm{MA}^{\mathrm{ext}}_\omega(u)\geq bf\,\mathrm{d}V$. For $p\in L^n$, the root functional is $\rho_\omega(p)=b_\omega(|p|^n)^{-1/n}$ if $p\ne0$, and $\rho_\omega(0)=0$. The space $\mathcal D_\omega$ of \cref{prop:individual-density-banach} consists of the $p\in L^n$ with $\rho_\omega(p)<+\infty$, and $\mathcal D_\omega^0$ is the closure of the bounded functions in $\mathcal D_\omega$, characterized by vanishing tails in \cref{prop:root-uniform-tails}. For a Gauduchon metric $\gamma$, $L_{\omega,\gamma}(p)$ denotes the test quotient of \cref{sec:integrable-gauduchon-tests}, and \cref{thm:integrable-gauduchon-formula} gives $\rho_\omega(p)=\sup_\gamma L_{\omega,\gamma}(|p|)$. Finally, $k_T$ is the determinant density of $T\in\mathscr G_\omega$ from \cref{thm:gaud-determinant-subcritical}.

If $p\geq0$ is measurable, define
\[
 D_p(T)=\int_Xp k_T\,\mathrm{d}V\qquad(T\in\mathscr G_\omega),
\]
allowing the value $+\infty$. If $T=\gamma^{n-1}/\int_X\omega\wedge\gamma^{n-1}$ for a Gauduchon metric $\gamma$, then $D_p(T)=L_{\omega,\gamma}(p)$.

\Needspace{11\baselineskip}
\begin{theorem}\label{thm:integrable-current-attainment}\index{Gauduchon duality!current attainment}
For every nonnegative $p\in L^n$,
\begin{equation}\label{eq:integrable-current-supremum}
 \sup_{T\in\mathscr G_\omega}D_p(T)=\rho_\omega(p),
\end{equation}
with the value $+\infty$ allowed. If $p\in\mathcal D_\omega^0$, the functional $D_p$ is finite, concave and weakly upper semicontinuous on $\mathscr G_\omega$, its maximum is attained, and the densities $p k_T$, for $T\in\mathscr G_\omega$, are uniformly integrable.
\end{theorem}
\begin{proof}
For a smooth strictly positive $p$, we have $D_p=\mathfrak D_{n\log p}$ and $\mathrm{e}^{-B_0(n\log p)/n}=\rho_\omega(p)$, so \cref{thm:weak-gaud-compact-duality}, applied to $F=n\log p$, proves \eqref{eq:integrable-current-supremum}. By H\"older's inequality and \cref{thm:gaud-determinant-subcritical}, for every real $r>n$ there is a constant $C_r$ such that, for every $q\in L^r$,
\begin{equation}\label{eq:current-root-test-bound}
 \sup_{T\in\mathscr G_\omega}\int_X|q| k_T\,\mathrm{d}V
       \leq C_r\|q\|_{L^r}.
\end{equation}
Next, approximate a bounded nonnegative $p$ in $L^{2n}$ by smooth strictly positive functions. This estimate and \eqref{eq:root-subcritical-bound} give uniform convergence on $\mathscr G_\omega$ of the corresponding functionals $D$ and convergence of the root functionals, so the supremum formula holds for bounded $p$. For such $p$, the functional $D_p$ is also upper semicontinuous and concave, by \cref{lem:weak-gaud-determinant-dual} and uniform convergence. For general $p\in L^n$, we use the truncations $\min(p,K)$ at positive integer levels $K$. Monotone convergence and the definition of $b_\omega$ through truncations give
\[
 \sup_T D_p(T)=\sup_T\sup_K D_{\min(p,K)}(T)
       =\sup_K\rho_\omega(\min(p,K))=\rho_\omega(p).
\]
This step only interchanges two suprema and gives no attainment.

Now suppose $p\in\mathcal D_\omega^0$. The supremum formula applied to the tail gives
\[
 0\leq D_p(T)-D_{\min(p,K)}(T)
 \leq \rho_\omega(p\mathbf1_{\{p>K\}})\longrightarrow0
\]
uniformly in $T$, by \cref{prop:root-uniform-tails}. Hence $D_p$ is finite and is a uniform limit of upper semicontinuous concave functionals, and compactness of $\mathscr G_\omega$ gives a maximizer. Finally, fix $1<q<n/(n-1)$. By H\"older's inequality and \cref{thm:gaud-determinant-subcritical}, for every measurable set $E$,
\[
 \sup_T\int_Ep k_T\,\mathrm{d}V
 \leq\rho_\omega(p\mathbf1_{\{p>K\}})
       +K C_q\,\mathrm{d}V(E)^{1-1/q}.
\]
Choosing first $K$ large and then the volume of $E$ small proves uniform absolute continuity. Since the total masses are bounded by $\rho_\omega(p)$, the densities are uniformly integrable.
\end{proof}

For a nonzero density, a maximizing current gives a supporting probability for the threshold, as we show below. Whether the threshold is attained by an exact equation depends on the regularity of the maximizing current, which enters through the trace inequalities of \cref{subsec:threshold-trace-inequalities}.

\subsection{The tail condition and perturbations of the density}

Vanishing tails are equivalent to continuity of the threshold under uniformly bounded multiplicative perturbations that converge in measure. This characterizes the precise role of the tail condition: if it fails, a perturbation supported on a shrinking set can lower the threshold by a fixed factor.

\begin{proposition}\label{prop:root-tail-continuity}
Let $0\leq p\in\mathcal D_\omega$ be nonzero. The following are equivalent:
\begin{enumerate}
\item $p\in\mathcal D_\omega^0$;
\item the densities $p k_T$, for $T\in\mathscr G_\omega$, are uniformly integrable;
\item whenever $F_j$ are smooth real functions, uniformly bounded and tending to zero in measure, one has
\[
 b_\omega(\mathrm{e}^{F_j}p^n)\longrightarrow b_\omega(p^n).
\]
\end{enumerate}
If the first condition holds, the convergence in the third also holds for bounded measurable $F_j$ with the same uniform bound and convergence in measure. If it fails, there are smooth $F_j\geq0$, uniformly bounded and tending to zero in measure, such that
\[
 \mathrm{e}^{F_j}p^n\longrightarrow p^n\quad\text{in }L^1,
 \qquad
 b_\omega(\mathrm{e}^{F_j}p^n)\leq2^{-n}b_\omega(p^n).
\]
\end{proposition}
\begin{proof}
The first condition implies the second by \cref{thm:integrable-current-attainment}. Conversely, $\mathrm{d}V(\{p>K\})\to0$, so uniform integrability and \eqref{eq:integrable-current-supremum} give
\[
 \rho_\omega(p\mathbf1_{\{p>K\}})
 =\sup_{T\in\mathscr G_\omega}\int_{\{p>K\}}p k_T\,\mathrm{d}V
 \longrightarrow0.
\]
This is the first condition by \cref{prop:root-uniform-tails}.

If $p\in\mathcal D_\omega^0$ and $|F_j|\leq A$, then $p(\mathrm{e}^{F_j/n}-1)$ tends to zero in measure and its absolute value is at most $(\mathrm{e}^{A/n}+1)p$. The dominated-convergence criterion of \cref{prop:root-uniform-tails} makes its root norm tend to zero. The triangle inequality gives $\rho_\omega(p\mathrm{e}^{F_j/n})\to\rho_\omega(p)>0$. Taking negative $n$th powers proves the third condition, including measurable perturbations.

Suppose now that the tails do not tend to zero, and put $R=\rho_\omega(p)>0$. The smooth Gauduchon formula yields $\delta>0$, numbers $K_j\to\infty$ and smooth strictly positive $T_j\in\mathscr G_\omega$ such that
\[
 \int_{\{p>K_j\}}p k_{T_j}\,\mathrm{d}V>\delta.
\]
The finite measure $p k_{T_j}\,\mathrm{d}V$ is regular. Choose a compact set $C_j\subset\{p>K_j\}$ whose mass exceeds $\delta/2$, and an open neighborhood $O_j$ with
$\mathrm{d}V(O_j)\leq\mathrm{d}V(\{p>K_j\})+1/j$. Take a smooth cutoff $0\leq\chi_j\leq1$ that equals one on $C_j$ and is supported in $O_j$. Choose a fixed $A>0$ with $\mathrm{e}^{A/n}\delta/2>2R$, and set $F_j=A\chi_j$. Then $F_j\to0$ in measure, while
\[
 \rho_\omega(p\mathrm{e}^{F_j/n})
 \geq\int_Xp\mathrm{e}^{F_j/n}k_{T_j}\,\mathrm{d}V
 >2R.
\]
The threshold formula gives the asserted uniform gap. Finally,
$|\mathrm{e}^{F_j}-1|p^n\leq(\mathrm{e}^A-1)p^n\mathbf1_{O_j}$, whose integral tends to zero by absolute continuity of the integral. This proves the last assertion and the converse to the third condition.
\end{proof}

For $f=p^n$, these perturbations satisfy $\mathrm{e}^{-A}f\leq\mathrm{e}^{F_j}f\leq\mathrm{e}^Af$ with a fixed $A$. The continuity statement is restricted to such an order interval. The unrestricted $L^1$ instability in \cref{sec:individual-density-stability} allows approximations without these bounds.

For attainment at one density, uniform integrability of a single maximizing sequence is enough. This weaker condition also identifies the limit of its weighted measures.

\begin{proposition}\label{prop:selected-gauduchon-attainment}
Let $0\leq p\in\mathcal D_\omega$ be nonzero and put $R=\rho_\omega(p)$. Suppose $T_j\in\mathscr G_\omega$ converge weakly to $T$ and satisfy $D_p(T_j)\to R$. If the densities $p k_{T_j}$ are uniformly integrable, then
\[
 D_p(T)=R,\qquad
 p k_{T_j}\,\mathrm{d}V\rightharpoonup p k_T\,\mathrm{d}V.
\]
Consequently, a uniformly integrable maximizing sequence has a maximizing current among its weak cluster limits.
\end{proposition}
\begin{proof}
For positive integers $K$, put $p_K=\min(p,K)$. Since $p$ is finite almost everywhere, uniform integrability gives
\[
 \varepsilon_K:=\sup_j\int_{\{p>K\}}p k_{T_j}\,\mathrm{d}V
       \longrightarrow0.
\]
For every nonnegative continuous $h$, the weight $hp_K$ is bounded and measurable. The bounded-weight upper semicontinuity proved in \cref{thm:integrable-current-attainment} therefore gives
\[
 \limsup_j\int_Xhp k_{T_j}\,\mathrm{d}V
 \leq \int_Xhp_K k_T\,\mathrm{d}V
                   +\|h\|_\infty\varepsilon_K.
\]
Let $K\to\infty$. For $h=1$ this yields $R\leq D_p(T)$; the supremum formula \eqref{eq:integrable-current-supremum} gives the reverse inequality. If $\mu$ is a weak cluster limit of $p k_{T_j}\,\mathrm{d}V$, the same estimate for all nonnegative continuous $h$ gives $\mu\leq p k_T\,\mathrm{d}V$. Both measures have mass $R$, so they agree. Their bounded masses give weak compactness, and every cluster limit has now been identified. The final assertion follows from compactness of $\mathscr G_\omega$.
\end{proof}


\subsection{Supporting probabilities for an integrable density}

Under the vanishing-tail condition, we can also perturb the density by continuous multiplicative factors $\mathrm{e}^F$. The logarithm of the threshold is then a concave function of $F$, and its supporting probabilities are given by the maximizing currents, just as for the continuous densities in \cref{sec:gaud-current-adjoints}.

Let $f\in L^1(\mathrm{d}V)$ be nonzero and nonnegative, and assume $p=f^{1/n}\in\mathcal D_\omega^0$. For a continuous real $F$, set
\[
 B_f(F)=\log b_\omega(\mathrm{e}^Ff)
       =-n\log\rho_\omega(p \mathrm{e}^{F/n}),\qquad
 M_f(F)=\rho_\omega(p \mathrm{e}^{F/n}).
\]
Both quantities are finite and $M_f(F)>0$, since $\mathcal D_\omega^0$ is an order ideal. Let $\mathscr M_f(F)$ be the set of currents in $\mathscr G_\omega$ maximizing $D_{p \mathrm{e}^{F/n}}$, which is nonempty by \cref{thm:integrable-current-attainment}, and put
\[
 \nu_{T,F}=\frac{p \mathrm{e}^{F/n} k_T}{M_f(F)}\,\mathrm{d}V
 \quad(T\in\mathscr M_f(F)).
\]
These are probability measures. Let $\mathcal A_f(F)$, the set of \emph{supporting probabilities} of $B_f$ at $F$, consist of all Borel probability measures $\nu$ satisfying
\begin{equation}\label{eq:integrable-threshold-support}
 B_f(G)\leq B_f(F)-\int_X(G-F)\,\mathrm{d}\nu
 \qquad(G\in C(X,\mathbb R)).
\end{equation}

\Needspace{20\baselineskip}
\begin{proposition}\label{prop:integrable-threshold-adjoints}
The functional $B_f$ is concave, decreasing and $1$-Lipschitz for the uniform norm, and $B_f(F+a)=B_f(F)-a$ for constants $a$. Moreover,
\[
 \mathcal A_f(F)=\{\nu_{T,F}:T\in\mathscr M_f(F)\}.
\]
This set is nonempty, convex and weakly compact, and its measures are absolutely continuous. For each $R>0$, the densities of all these measures, for all $F$ with $\|F\|_\infty\leq R$, are uniformly integrable. For $H\in C(X,\mathbb R)$, the one-sided directional derivatives of $B_f$ at $F$ are
\[
 (B_f)'_+(F;H)=-\max_{\nu\in\mathcal A_f(F)}\int_XH\,\mathrm{d}\nu,
 \qquad
 (B_f)'_-(F;H)=-\min_{\nu\in\mathcal A_f(F)}\int_XH\,\mathrm{d}\nu.
\]
\end{proposition}
\begin{proof}
For each $T$, H\"older's inequality makes $F\mapsto\int_Xp \mathrm{e}^{F/n}k_T\,\mathrm{d}V$ logarithmically convex, and taking suprema preserves the inequality defining logarithmic convexity. Thus $M_f$ is logarithmically convex and $B_f=-n\log M_f$ is concave. Monotonicity and homogeneity of $\rho_\omega$ give the translation rule and, by comparison with constant translates, the Lipschitz estimate.

For fixed $F$, the weight $p\mathrm{e}^{F/n}$ belongs to $\mathcal D_\omega^0$, so the densities $p\mathrm{e}^{F/n}k_T$, $T\in\mathscr G_\omega$, are uniformly integrable by \cref{thm:integrable-current-attainment}. If $T_j\to T$ and $D_{p\mathrm{e}^{F/n}}(T_j)\to M_f(F)$, \cref{prop:selected-gauduchon-attainment} makes $T$ a maximizer and gives convergence of the weighted measures. For $T_1,T_2\in\mathscr G_\omega$ and real $0\leq t\leq1$, matrix determinant concavity gives
\[
 p \mathrm{e}^{F/n}k_{tT_1+(1-t)T_2}\,\mathrm{d}V
 \geq t p \mathrm{e}^{F/n}k_{T_1}\,\mathrm{d}V
       +(1-t)p \mathrm{e}^{F/n}k_{T_2}\,\mathrm{d}V.
\]
For two maximizing currents, the current on the left also maximizes, and equal total masses turn this inequality into an equality. The set $\mathscr N_f(F)=\{\nu_{T,F}:T\in\mathscr M_f(F)\}$ is therefore nonempty, compact and convex.

\Needspace{6\baselineskip}
A maximizing $T$ gives
\[
 M_f(G)\geq M_f(F)\int_X\mathrm{e}^{(G-F)/n}\,\mathrm{d}\nu_{T,F}.
\]
Taking $-n$ times the logarithm and applying Jensen's inequality proves \eqref{eq:integrable-threshold-support}, so $\mathscr N_f(F)\subset\mathcal A_f(F)$. For the reverse inclusion, a Taylor expansion of the exponential gives
\[
 D_{p \mathrm{e}^{(F+tH)/n}}(T)
 =D_{p \mathrm{e}^{F/n}}(T)+\frac tn\int_X H p \mathrm{e}^{F/n} k_T\,\mathrm{d}V+O(t^2),
\]
where the error is uniform in $T$ since $D_{p \mathrm{e}^{F/n}}(T)\leq M_f(F)$. Maximizers at $F+tH$ are near-maximizers at $F$, by the multiplicative bounds $\mathrm{e}^{-t\|H\|_\infty/n}\leq M_f(F+tH)/M_f(F)\leq\mathrm{e}^{t\|H\|_\infty/n}$. The convergence of their weighted measures proved above and this uniform expansion are exactly the two inputs of the one-sided differentiation argument in the proof of \cref{thm:gaud-current-all-adjoints}, which gives
\[
 (B_f)'_+(F;H)=-\max_{\nu\in\mathscr N_f(F)}\int_XH\,\mathrm{d}\nu.
\]
The separation argument of that proof then identifies all supporting probabilities with $\mathscr N_f(F)$, and replacing $H$ by $-H$ gives the left derivative.

Finally, $M_f(F)\geq \mathrm{e}^{-R/n}\rho_\omega(p)$ when $\|F\|_\infty\leq R$, so the densities of the probabilities $\nu_{T,F}$ are at most $\mathrm{e}^{2R/n}\rho_\omega(p)^{-1}p k_T$. The uniform integrability of the densities $pk_T$ in \cref{thm:integrable-current-attainment} proves the remaining assertions.
\end{proof}

When $f=1$, we have $B_f=B_0$, and the proposition reduces to \cref{thm:gaud-current-all-adjoints}. For general $f$, the functional $B_f$ is defined through the subsolution threshold and does not presuppose that the corresponding equation has a solution. If $v_-(\omega)>0$, every nonzero nonnegative $f\in L^1$ satisfies $f^{1/n}\in\mathcal D_\omega^0$ by \cref{prop:root-uniform-tails}, and \cref{cor:positive-lower-mass-exact-threshold} identifies $b_\omega(f)$ as the unique constant for which the equation $\mathrm{MA}^{\mathrm{ext}}_\omega(v)=b\,f\,\mathrm{d}V$ has a solution $v$ in the envelope class $\mathcal E_{\mathrm{env}}(X,\omega,0)$ of that corollary. Without positive lower mass, the uniform tail condition is an explicit sufficient hypothesis for this compact dual description.

\Needspace{16\baselineskip}
\subsection{Trace inequalities at the subsolution threshold}
\label{subsec:threshold-trace-inequalities}

For a nonzero $0\leq f\in L^1(\mathrm{d}V)$ with $b=b_\omega(f)>0$, recall from \cref{sec:individual-density-subsolutions} that a threshold subsolution is a function $u\in\PSH(X,\omega)$ with $\mathrm{MA}^{\mathrm{ext}}_\omega(u)\geq bf\,\mathrm{d}V$, and a threshold solution is one for which equality holds; threshold subsolutions exist by \cref{thm:individual-density-subsolution}. A smooth Gauduchon test can be paired with the full current $\omega+\ddc u$ of a threshold subsolution, and by \cref{lem:unbounded-gauduchon-trace-density} the resulting trace measure is bounded below by the determinant-weighted density. The nonnegative difference has a mass that depends only on the test. Along a maximizing sequence this mass tends to zero, uniformly over all threshold subsolutions. We write $\|\cdot\|_{\mathrm{TV}}$ for the total variation of a finite signed measure.

\Needspace{14\baselineskip}
\begin{proposition}\label{prop:threshold-trace-saturation}\index{Gauduchon duality!trace saturation}
Let $0\leq f\in L^1$ be nonzero with $b=b_\omega(f)>0$, write $p=f^{1/n}$, and let $u$ be a threshold subsolution. For every smooth strictly positive $T\in\mathscr G_\omega$, the measure
\[
 R_{u,T}=(\omega+\ddc u)\wedge T-b^{1/n}p k_T\,\mathrm{d}V
\]
is positive and has mass
\[
 R_{u,T}(X)=1-b^{1/n}D_p(T).
\]
In particular, if $(T_j)$ is a sequence of smooth strictly positive members of $\mathscr G_\omega$ with $D_p(T_j)\to\rho_\omega(p)$, then $\|R_{u,T_j}\|_{\mathrm{TV}}\to0$, uniformly over all threshold subsolutions $u$.

If also $p\in\mathcal D_\omega^0$ and $T_j\to T$ as currents, then $T$ maximizes $D_p$ and
\[
 (\omega+\ddc u)\wedge T_j
       \rightharpoonup b^{1/n}p k_T\,\mathrm{d}V.
\]
\end{proposition}
\begin{proof}
A smooth strictly positive $T\in\mathscr G_\omega$ is $\gamma^{n-1}$ for a Gauduchon metric $\gamma$ normalized by $\int_X\omega\wedge\gamma^{n-1}=1$. \Cref{lem:unbounded-gauduchon-trace-density} proves positivity of $R_{u,T}$, and Stokes' theorem for currents gives $((\omega+\ddc u)\wedge T)(X)=\int_X\omega\wedge T=1$. This proves the mass formula. Since $\rho_\omega(p)=b^{-1/n}$, the nonnegative mass tends to zero along a maximizing sequence; it equals the total variation norm and does not depend on $u$.

Under the tail hypothesis, \cref{thm:integrable-current-attainment} gives uniform integrability of the densities $p k_{T_j}$. Thus \cref{prop:selected-gauduchon-attainment} makes $T$ a maximizer and gives
$p k_{T_j}\,\mathrm{d}V\rightharpoonup p k_T\,\mathrm{d}V$. Adding the error measures $R_{u,T_j}$, which tend to zero in total variation, proves the limit formula.
\end{proof}

The limit formula concerns smooth maximizing sequences. It does not define the product $(\omega+\ddc u)\wedge T$ for a singular current $T$, nor does it show that the limit is the same for different maximizing currents.

If a maximizing test is itself a smooth strictly positive form, the trace inequality becomes an equality of currents, and since $n\geq2$ this forces an exact smooth equation. The local regularity needed for this comes from \cref{lem:smooth-conformal-current}.

\Needspace{12\baselineskip}
\begin{corollary}\label{cor:smooth-dual-exact-equation}\index{Gauduchon duality!smooth maximizing test}\label{cor:smooth-dual-threshold-rigidity}
Let $0\leq f\in L^1(\mathrm{d}V)$ be nonzero with $b=b_\omega(f)>0$. If $D_{f^{1/n}}$ has a smooth strictly positive maximizer $T\in\mathscr G_\omega$, then $f^{1/n}$ agrees almost everywhere with a smooth nonnegative function, and every threshold subsolution $u$ is smooth and satisfies
\[
 \mathrm{MA}_\omega(u)=bf\,\mathrm{d}V.
\]
In particular a smooth threshold solution exists, and any two threshold subsolutions differ by a constant.
\end{corollary}
\Needspace{8\baselineskip}
\begin{proof}
Let $u$ be any threshold subsolution; one exists by \cref{thm:individual-density-subsolution}. Write $T=\gamma^{n-1}$, where $\gamma$ is the smooth Gauduchon metric given by the positive $(n-1)$st root of $T$, and set
\[
 a=\left(\frac{bf\,\mathrm{d}V}{\gamma^n}\right)^{1/n}.
\]
Put $A=\omega+\ddc u\geq0$. Since $T$ maximizes, the positive measure $R_{u,T}$ of \cref{prop:threshold-trace-saturation} has mass zero, which gives the equality
\[
 A\wedge\gamma^{n-1}=a\gamma^n.
\]
The coefficient measures of the positive current $A$ are controlled by this trace measure, which is absolutely continuous. Thus $A$ has a matrix density that is positive semidefinite almost everywhere, with coefficients in $L^n$ because $a\in L^n$.

\Needspace{5\baselineskip}
For every smooth Hermitian form $\beta$, \cref{lem:unbounded-gauduchon-trace-density} also gives
\[
 A\wedge\beta^{n-1}
 \geq\left(\frac{bf\,\mathrm{d}V}{\beta^n}\right)^{1/n}\beta^n.
\]
The resulting pointwise matrix inequalities hold almost everywhere, simultaneously for all positive definite matrices: on each coordinate chart, first use a countable dense family of constant positive matrices, realized on smaller charts by global smooth Hermitian forms, and then use continuity in the matrix. The variational formula for $(\det P)^{1/n}$ in the proof of \cref{lem:weak-gaud-determinant-dual} therefore yields
\[
 \left(\frac{A^n}{\gamma^n}\right)^{1/n}\geq a
 \quad\text{almost everywhere}.
\]
Here $A^n$ denotes only the determinant of the measurable coefficient matrix of $A$; no nonlinear product of currents is used. The trace equality gives the reverse inequality, by the arithmetic--geometric mean inequality. Equality then forces all eigenvalues of $A$ relative to $\gamma$ to equal $a$, also where $a=0$. Consequently $A=a\gamma$ as currents.

By \cref{lem:smooth-conformal-current}, applied with $\theta=\omega$, both $a$ and $u$ are smooth, and $f^{1/n}=b^{-1/n}a(\gamma^n/\mathrm{d}V)^{1/n}$ has a smooth nonnegative representative. Since $u$ is smooth, we may now take the top power: $(\omega+\ddc u)^n=a^n\gamma^n=bf\,\mathrm{d}V$. Finally, the function $a$ was defined from $f$, $b$ and $\gamma$ alone, so two threshold subsolutions $u_1,u_2$ satisfy $\omega+\ddc u_1=a\gamma=\omega+\ddc u_2$. Their smooth difference is therefore pluriharmonic, hence constant on the connected compact manifold $X$.
\end{proof}

If $f$ is smooth and strictly positive, \cref{thm:gaud-volume-dual-formula} provides a smooth strictly positive maximizing Gauduchon form, obtained from the smooth equation. The corollary therefore applies, and shows in particular that all threshold subsolutions, bounded or not, differ from the smooth solution by constants.


\begin{corollary}\label{cor:smooth-dual-exponential-limit}
Under the hypotheses of \cref{cor:smooth-dual-exact-equation}, let $u$ be the threshold solution with $\sup_Xu=0$. If $u_\lambda$ is the unique bounded solution of
\[
 \mathrm{MA}_\omega(u_\lambda)=\mathrm{e}^{\lambda u_\lambda}f\,\mathrm{d}V,
 \qquad\lambda>0,
\]
then $u_\lambda-\sup_Xu_\lambda\to u$ uniformly as $\lambda\downarrow0$.
\end{corollary}
\begin{proof}
By \cref{cor:smooth-dual-exact-equation}, $f$ agrees almost everywhere with a smooth function, so it belongs to $L^p$ for every finite $p$. By \cref{prop:individual-normalizer-lp}, the constant $c$ for which $\mathrm{MA}_\omega(v)=cf\,\mathrm{d}V$ has a bounded solution is $b_\omega(f)$. Thus \cref{prop:symmetric-equation-clusters}, with the trivial symmetry group, shows that the normalized family $u_\lambda-\sup_Xu_\lambda$ is uniformly precompact as $\lambda\downarrow0$ and that every cluster limit is a normalized solution with this constant. Each cluster limit is therefore a normalized threshold subsolution and equals $u$ by \cref{cor:smooth-dual-exact-equation}. If the whole family did not converge uniformly, a sequence staying at a fixed positive uniform distance from $u$ would have a uniformly convergent subsequence, a contradiction.
\end{proof}

\begin{corollary}\label{cor:kahler-smooth-dual-attainment}
Suppose $\omega$ is K\"ahler. For a nonzero nonnegative $f\in L^1(\mathrm{d}V)$, the functional $D_{f^{1/n}}$ has a smooth strictly positive maximizer in $\mathscr G_\omega$ if and only if $f$ agrees almost everywhere with a smooth strictly positive function.
\end{corollary}
\begin{proof}
If a smooth strictly positive maximizer exists, the proof of \cref{cor:smooth-dual-exact-equation} gives
\[
 \omega+\ddc u=a\gamma,
 \qquad a\geq0,
\]
with $a,u$ smooth and $a\not\equiv0$. Since $\mathrm{d}\omega=0$, the first-order equation in \cref{lem:smooth-conformal-current} becomes $\mathrm{d}a=-aC$ for a smooth one-form $C$. Along any smooth path $\sigma$, this is the linear scalar equation
\[
 (a\circ\sigma)'=-(a\circ\sigma)C(\dot\sigma),
\]
whose solution is the initial value multiplied by a positive exponential. Since $X$ is connected and $a\not\equiv0$, it follows that $a>0$ everywhere, and $f=b_\omega(f)^{-1}a^n\gamma^n/\mathrm{d}V$ is smooth and strictly positive. Conversely, if $f$ is smooth and strictly positive, \cref{thm:weak-gaud-compact-duality}, applied to $F=\log f$, for which $\mathfrak D_F=D_{f^{1/n}}$, provides a smooth strictly positive maximizer.
\end{proof}

For nonclosed backgrounds, the strict positivity of $f$ in \cref{cor:kahler-smooth-dual-attainment} can fail on every compact complex manifold of dimension at least two, including K\"ahler manifolds.

\begin{example}\label{ex:nonclosed-dual-zero-density}
On every connected compact complex manifold $X$ of dimension $n\geq2$, with a fixed smooth strictly positive probability volume $\mathrm{d}V$, there is a nonclosed Hermitian form $\omega$ and a nonzero smooth nonnegative density $f$ vanishing on an open set such that $b_\omega(f)=1$ and $D_{f^{1/n}}$ has a smooth strictly positive maximizer.
\end{example}
\begin{proof}
Fix a Gauduchon metric $\gamma$ and a point $x$. Choose a smooth real function $\psi$ such that $\ddc\psi\leq-\gamma$ on a coordinate neighborhood $U$ of $x$; in holomorphic coordinates $z=(z_1,\ldots,z_n)$ centered at $x$, such a function is obtained by multiplying $-M|z|^2$, with $M$ large enough, by a cutoff that equals one near $x$, and then shrinking $U$.

Choose a smooth function $\chi$ with $0\leq\chi\leq1$ that vanishes on a smaller neighborhood of $x$ and equals one outside $U$. Put
\[
 a=L\chi,\qquad \omega=a\gamma-\ddc\psi.
\]
Since $X$ is compact, $\omega$ is strictly positive for sufficiently large $L>0$: on $U$ the term $-\ddc\psi$ is at least $\gamma$, and outside $U$ the term $L\gamma$ dominates $\ddc\psi$. Set $f=a^n\gamma^n/\mathrm{d}V$. The smooth function $\psi$ satisfies $\omega+\ddc\psi=a\gamma\geq0$ and $\mathrm{MA}_\omega(\psi)=f\,\mathrm{d}V$, so $b_\omega(f)\geq1$.

Stokes' theorem and $\ddc\gamma^{n-1}=0$ give
\[
 \int_X\omega\wedge\gamma^{n-1}
 =\int_Xa\gamma^n
 =\int_X(f\,\mathrm{d}V/\gamma^n)^{1/n}\gamma^n.
\]
Testing \eqref{eq:integrable-gauduchon-threshold} with $\gamma$ therefore gives the opposite inequality $b_\omega(f)\leq1$. Hence $b_\omega(f)=1$, the metric $\gamma$ attains the infimum, and the normalized current $\gamma^{n-1}/\int_X\omega\wedge\gamma^{n-1}$ is a smooth strictly positive maximizer of $D_{f^{1/n}}$. If $\omega$ were closed, \cref{cor:kahler-smooth-dual-attainment} would force $f$ to be strictly positive, contrary to its vanishing on an open set. Thus $\omega$ is nonclosed.
\end{proof}

In \cref{ex:nonclosed-dual-zero-density}, \cref{cor:smooth-dual-exact-equation,cor:smooth-dual-exponential-limit} show that every threshold subsolution is a constant translate of $\psi$ and that the normalized solutions of the exponential equations converge uniformly to $\psi-\sup_X\psi$. Thus vanishing of the density on an open set does not by itself force nonuniqueness or the failure of this approximation.

The restriction $n\geq2$ in \cref{cor:smooth-dual-exact-equation} is necessary. On a compact connected Riemann surface $X$ with a Hermitian form $\omega$, a normalized Gauduchon current has bidegree $(0,0)$, and the equation $\ddc T=0$ makes it a harmonic function, hence the constant $(\int_X\omega)^{-1}$. This single smooth strictly positive current maximizes for every nonzero nonnegative $L^1$ density, which need not have a smooth representative.


\subsection{A smooth maximizing form with an analytic degeneracy set}

Strict positivity of a smooth maximizing form can be weakened if its degeneracy is confined to a proper analytic subset. The proof is local away from that subset. The possible ambiguity of the potential is then divisorial.

\begin{theorem}\label{thm:semipositive-dual-exact-equation}
Let $\omega$ be Hermitian on a connected compact complex manifold of dimension $n\geq2$, let $0\leq f\in L^1(\mathrm{d}V)$ be nonzero, and assume $b=b_\omega(f)>0$. Suppose $D_{f^{1/n}}$ has a smooth semipositive maximizer $T\in\mathscr G_\omega$ which is strictly positive on $X\setminus Z$, where $Z$ is a proper analytic subset. Then every threshold subsolution is smooth on $X\setminus Z$ and satisfies
\[
 \mathrm{MA}^{\mathrm{ext}}_\omega(u)=bf\,\mathrm{d}V
\]
on $X$. Such a solution exists and has finite mass. The function $f^{1/n}$ has a smooth nonnegative representative on $X\setminus Z$.

If $u,v$ are threshold subsolutions and $Z_1,\ldots,Z_m$ are the prime divisorial components of $Z$, there are real numbers $c_j$ such that
\begin{equation}\label{eq:semipositive-dual-divisorial-difference}
 \ddc(u-v)=\sum_{j=1}^m c_j[Z_j],\qquad
 \sum_{j=1}^m c_j\{Z_j\}=0
       \quad\text{in }H^{1,1}_{\mathrm{BC}}(X,\mathbb R).
\end{equation}
Consequently threshold solutions are unique modulo constants within each multiplier-ideal equivalence class, hence within each ordinary singularity type. They are unique modulo constants without a prescribed type if the classes $\{Z_j\}$ are linearly independent. In particular this holds when $Z$ has codimension at least two.
\end{theorem}
\begin{proof}
Put $p=f^{1/n}$ and $A=\omega+\ddc u$, where $u$ is any threshold subsolution. Existence of at least one such $u$ is \cref{thm:individual-density-subsolution}. Fix a smooth strictly positive $T_0\in\mathscr G_\omega$ and set $T_\varepsilon=(T+\varepsilon T_0)/(1+\varepsilon)$. Each $T_\varepsilon$ belongs to $\mathscr G_\omega$ and has a smooth Hermitian root, so \cref{prop:threshold-trace-saturation} gives
\[
 A\wedge T_\varepsilon\geq b^{1/n}p k_{T_\varepsilon}\,\mathrm{d}V.
\]
For smooth semipositive forms, $k_T$ is the determinant root of the coefficient matrix, in the normalization of \cref{sec:weak-gauduchon-duality}. This is a continuous function of that matrix. Thus $k_{T_\varepsilon}\to k_T$ uniformly on compact $X$. Since $p\in L^1$, passage to the limit gives
\[
 A\wedge T\geq b^{1/n}p k_T\,\mathrm{d}V.
\]
The left side has mass one by Stokes and $\ddc T=0$. The right side has mass $b^{1/n}D_p(T)=1$, because $T$ maximizes. Their difference is positive and has zero mass, hence
\begin{equation}\label{eq:semipositive-dual-trace-equality}
 A\wedge T=b^{1/n}p k_T\,\mathrm{d}V.
\end{equation}

On $U=X\setminus Z$ write $T=\gamma^{n-1}$, where $\gamma$ is smooth and strictly positive, and set $a=(bf\,\mathrm{d}V/\gamma^n)^{1/n}$. The trace equality becomes $A\wedge\gamma^{n-1}=a\gamma^n$. On relatively compact subsets of $U$ this controls every coefficient measure of $A$ by an $L^n$ density. The all-Hermitian trace inequalities and matrix arithmetic--geometric mean, exactly as in the proof of \cref{cor:smooth-dual-exact-equation}, force $A=a\gamma$ there. This step uses only the local positivity of $\gamma$; the other Hermitian tests extend smoothly from smaller coordinate charts to $X$.

The local regularity lemma \cref{lem:smooth-conformal-current} makes $u,a$ smooth on $U$. Taking the smooth top power proves the equation there, and $f^{1/n}=b^{-1/n}a(\gamma^n/\mathrm{d}V)^{1/n}$ is smooth and nonnegative there. Both $\mathrm{MA}^{\mathrm{ext}}_\omega(u)$ and $f\,\mathrm{d}V$ vanish on the pluripolar set $Z$. Hence the equality holds on $X$, and its total mass is $b\int_Xf\,\mathrm{d}V<+\infty$.

For two threshold subsolutions the formula $\omega+\ddc u=a\gamma$ has the same right side on $U$. Thus $S=\ddc(u-v)$ is a real closed current supported on $Z$. It has order zero, being the difference of the positive currents $\omega+\ddc u$ and $\omega+\ddc v$. The classical support theorem gives $S=\sum_jc_j[Z_j]$; components of higher codimension carry no such current. This is \cite[Chapter III, Corollaries 2.11 and 2.14]{DemCADG}. To check the class relation with smooth representatives, write $[Z_j]=\beta_j+\ddc\ell_j$, using a smooth Hermitian metric on the divisor line bundle. Then
\[
 \sum_jc_j\beta_j
   =\ddc\bigl(u-v-\sum_jc_j\ell_j\bigr).
\]
The left side is smooth. Taking its local trace and applying elliptic regularity makes the global $L^1$ potential on the right smooth. Hence the displayed form is $\ddc$-exact in smooth Bott--Chern cohomology, proving \eqref{eq:semipositive-dual-divisorial-difference}.

Near a regular point of $Z_j$ lying on no other component, choose a local defining function $z_1$. In our normalization $\ddc\log|z_1|^2=[Z_j]$, so
\[
 \ddc\bigl(u-v-c_j\log|z_1|^2\bigr)=0.
\]
The expression in parentheses is a smooth pluriharmonic function by distributional elliptic regularity. Taking generic Lelong numbers therefore gives $c_j=\nu_{Z_j}(u)-\nu_{Z_j}(v)$. If $u\sim_{\mathcal I}v$, these numbers agree by \cref{thm:local-i-inputs}, so all coefficients vanish. Ordinary equivalence implies multiplier-ideal equivalence. Alternatively, independence of the divisor classes forces them all to vanish directly from \eqref{eq:semipositive-dual-divisorial-difference}. In either case $\ddc(u-v)=0$ globally, and a real pluriharmonic distribution on connected compact $X$ is constant. When $Z$ has codimension at least two, the sum is empty.
\end{proof}

The theorem assumes a smooth form. It neither defines a wedge product with an arbitrary singular maximizing current nor proves that such a current has an analytic degeneracy set. The divisor relation is necessary for two threshold solutions to differ; it does not assert that every relation is realized. No bounded mass, positive lower mass, root-tail condition, or convergence of an exponential approximation is used.


\Needspace{22\baselineskip}
\subsection{The limit of the exponential equations}
\label{subsec:semipositive-exponential-setting}

Use the data of \cref{thm:semipositive-dual-exact-equation}: the manifold $X$ of dimension $n\geq2$, Hermitian form $\omega$, probability volume $\mathrm{d}V$, nonzero nonnegative density $f$, positive threshold $b=b_\omega(f)$, and smooth semipositive maximizing form $T$ with analytic degeneracy set $Z$. Put $p=f^{1/n}$ and let $k_T$ be its determinant density. All $L^1$ convergence below uses $\mathrm{d}V$.

A smooth semipositive maximizing form also identifies the normalized exponential limit. The argument passes the linear trace inequalities to the limit before using the local regularity proof.

\begin{proposition}\label{prop:semipositive-dual-exponential-limit}
Assume the hypotheses of \cref{thm:semipositive-dual-exact-equation}, and let $u_0$ be its unique threshold solution with $\sup_Xu_0=0$. For real $\lambda>0$, let $u_\lambda$ be the greatest exponential solution of \cref{cor:individual-positive-threshold-exponential}, and put
\[
 v_\lambda=u_\lambda-\sup_Xu_\lambda,\qquad
 c_\lambda=\mathrm{e}^{\lambda\sup_Xu_\lambda}.
\]
There is a constant $C$ independent of $\lambda>0$ such that
\[
 b\leq c_\lambda\leq b\,\mathrm{e}^{C\lambda},\qquad
 u_0-C\leq v_\lambda\leq0.
\]
Moreover $v_\lambda\to u_0$ in $L^1(X)$ as $\lambda\downarrow0$.
\end{proposition}
\begin{proof}
Write $p=f^{1/n}$ and introduce the probability measure
\[
 \mathrm{d}\nu=b^{1/n}p k_T\,\mathrm{d}V.
\]
It has mass one because $T$ maximizes. Its density belongs to $L^n$, since $p\in L^n$ and the smooth semipositive form $T$ has bounded determinant density $k_T$. By \cref{lem:uniform-exponential}, normalized $\omega$-plurisubharmonic functions have uniformly bounded moments of every finite order. H\"older's inequality with exponents $n$ and $n/(n-1)$ therefore gives
\[
 \int_X(-v_\lambda)\,\mathrm{d}\nu\leq C.
\]
Here $n\geq2$ is used.

Apply the trace inequality to $u_\lambda$, or equivalently to $v_\lambda$, and pass from positive regularizations of $T$ to $T$ as in the preceding proof. After integration,
\[
 1\geq(c_\lambda/b)^{1/n}
       \int_X\mathrm{e}^{\lambda v_\lambda/n}\,\mathrm{d}\nu
 \geq(c_\lambda/b)^{1/n}\mathrm{e}^{-C\lambda/n},
\]
where the last step is Jensen's inequality. This proves the upper scalar bound. The lower bound and the comparison
\[
 v_\lambda\geq u_0-\lambda^{-1}\log(c_\lambda/b)
\]
are \cref{cor:individual-positive-threshold-exponential}, giving the asserted common lower barrier.

Take any sequence $\lambda_j\downarrow0$ and an $L^1$-convergent subsequence $v_{\lambda_j}\to v$, supplied by \cref{thm:normalized-compactness}. Its limit has supremum zero. The normalized $L^1$ bound implies $\lambda v_\lambda\to0$ in measure. Since $v_\lambda\leq0$ and $p\in L^1$, it follows that
\[
 c_\lambda^{1/n}\mathrm{e}^{\lambda v_\lambda/n}p
       \longrightarrow b^{1/n}p
       \quad\text{in }L^1(\mathrm{d}V).
\]
Put $A_\lambda=\omega+\ddc v_\lambda$. The positive trace errors
\[
 A_\lambda\wedge T
   -c_\lambda^{1/n}\mathrm{e}^{\lambda v_\lambda/n}p k_T\,\mathrm{d}V
\]
have masses tending to zero, since the two terms have total masses tending to one. The $L^1$ convergence of the potentials gives weak convergence of their linear currents. Consequently
\[
 (\omega+\ddc v)\wedge T=b^{1/n}p k_T\,\mathrm{d}V.
\]
For every smooth Hermitian form $\beta$, the all-Hermitian trace inequality passes to the same limit and gives
\[
 (\omega+\ddc v)\wedge\beta^{n-1}
   \geq(bf\,\mathrm{d}V/\beta^n)^{1/n}\beta^n.
\]
The local matrix and regularity argument in the proof of \cref{thm:semipositive-dual-exact-equation} now applies on $X\setminus Z$: with $\gamma$ the positive root of $T$ there and $a=(bf\,\mathrm{d}V/\gamma^n)^{1/n}$, it gives $\omega+\ddc v=a\gamma$, then smoothness, and then the exact equation there. Both extended measures ignore $Z$, so $v$ is a global threshold solution. Uniqueness and normalization give $v=u_0$.

Every sequence tending to zero has a subsequence with this same $L^1$ limit. Normalized compactness therefore proves convergence of the whole family.
\end{proof}

The lower barrier gives capacity convergence under the additional envelope condition below. If the maximizing form is strictly positive everywhere, \cref{cor:smooth-dual-exponential-limit} already gives uniform convergence.


\Needspace{15\baselineskip}
\begin{corollary}\label{cor:semipositive-dual-capacity-limit}
Under the hypotheses of \cref{prop:semipositive-dual-exponential-limit}, suppose in addition that $P_{\delta\omega}(u_0)$ is genuine for some $0<\delta<1$. Then $v_\lambda\to u_0$ in capacity as $\lambda\downarrow0$. In particular this holds if the threshold solution $u_0$ is bounded.
\end{corollary}
\begin{proof}
Fix a sequence $\lambda_j\downarrow0$. \Cref{prop:semipositive-dual-exponential-limit} gives $u_0-C\leq v_{\lambda_j}\leq0$ and $v_{\lambda_j}\to u_0$ in $L^1$. For all large $j$,
\[
 \mathrm{MA}^{\mathrm{ext}}_\omega(v_{\lambda_j})
   =c_{\lambda_j}\mathrm{e}^{\lambda_jv_{\lambda_j}}f\,\mathrm{d}V
   \leq b\mathrm{e}^{C}f\,\mathrm{d}V.
\]
The bounded functions $|\mathrm{e}^{v_{\lambda_j}}-\mathrm{e}^{u_0}|$ tend to zero in smooth-volume measure. Since $f\in L^1$, their integrals against $f\,\mathrm{d}V$ tend to zero. Thus the self-measure error in \eqref{eq:common-barrier-self-error} vanishes. Apply \cref{prop:lower-barrier-criterion} with $\theta=\omega$ and $\varphi=u_0$; its additional limit-type hypothesis is automatic because the limit is $u_0$ itself. This proves capacity convergence along every sequence and hence for the family. If $u_0$ is bounded, a constant below it is a competitor for $P_{\delta\omega}(u_0)$, so that envelope is genuine.
\end{proof}

The smaller-background condition is an additional hypothesis. Its validity for an arbitrary threshold solution in this semipositive dual setting is not established here.


\subsection{Envelope criteria and extension across the degeneracy set}
\label{subsec:semipositive-envelope-extension}

The capacity criterion above asks for a genuine smaller-background obstacle envelope. Two sufficient conditions come from different sources: invariant bounded mass controls the contact measures of stretched truncations, while a sublinear singularity bound supplies explicit competitors. Both conditions give zero Lelong numbers. Under dual attainment, a globally smooth root of rank at least two then permits smooth extension across the analytic degeneracy set.

We retain the connected compact manifold $X$ of dimension $n\geq2$, the Hermitian background $\omega$ and the smooth positive probability volume $\mathrm{d}V$. All obstacle inequalities concern canonical potentials. We will use one consequence of \cref{lem:gauduchon-uniform-lelong}: if, for every $0<\delta<1$, a genuine $v_\delta\in\PSH(X,\delta\omega)$ satisfies $v_\delta\leq u$, then
\[
 0\leq\nu(u,x)\leq\nu(v_\delta,x)
       =\delta\nu(v_\delta/\delta,x)\leq\delta C_\omega.
\]
Thus $u$ has zero Lelong numbers. The bound is uniform in $\delta$ because $v_\delta/\delta$ has the fixed background $\omega$.

\begin{theorem}\label{thm:invariant-mass-threshold-envelopes}\index{Envelope!invariant mass criterion}
Suppose $\ddc(\omega^j)=0$ for $1\leq j<n$. Put $V=\int_X\omega^n$. For every nonzero $0\leq f\in L^1(\mathrm{d}V)$, with $M=\int_Xf\,\mathrm{d}V$, one has
\[
 b_\omega(f)=\frac VM.
\]
Every threshold subsolution satisfies $\mathrm{MA}^{\mathrm{ext}}_\omega(u)=(V/M)f\,\mathrm{d}V$ and has mass $V$. Its obstacle envelope $P_{\delta\omega}(u)$ is genuine for every $0<\delta<1$, and $u$ has zero Lelong numbers. More generally, these last two conclusions hold for every $\omega$-plurisubharmonic potential of extended mass $V$.
\end{theorem}
\begin{proof}
By \cref{thm:mass-invariance-characterization}, every bounded $\omega$-plurisubharmonic potential has mass $V$. Truncation therefore bounds every extended mass above by $V$. For $f_k=\min(f,k)$, integration of the bounded equation gives $c_\omega(f_k)=V/\int_Xf_k\,\mathrm{d}V$. The definition \eqref{eq:individual-truncated-normalizer} gives $b_\omega(f)=V/M$. A threshold subsolution has mass at least $b_\omega(f)M=V$. Its mass is consequently $V$, and the nonnegative difference from $b_\omega(f)f\,\mathrm{d}V$ has zero mass.

Let now $u$ have extended mass $V$. Normalize $u\leq0$, put $\mu=\mathrm{MA}^{\mathrm{ext}}_\omega(u)$, and fix $a>1$. For $u_j=\max(u,-j)$, let $w_j=P_\omega(au_j)$. Constants give $-aj\leq w_j\leq0$. The bounded contact theorem \cref{thm:bounded-contact-support} concentrates its measure on $C_j=\{w_j=au_j\}$. Canonical almost-everywhere order gives $w_j\leq au_j$ everywhere. Both potentials are $a\omega$-plurisubharmonic, so positivity and ordered contact give
\[
 \mathrm{MA}_\omega(w_j)
 \leq\mathbf1_{C_j}\mathrm{MA}_{a\omega}(w_j)
 \leq a^n\mathbf1_{C_j}\mathrm{MA}_\omega(u_j).
\]
For integers $j\geq N$, locality and the bounded mass identity yield
\[
 \int_{\{u_j\leq-N\}}\mathrm{MA}_\omega(u_j)
       =V-\mu(\{u>-N\})=\mu(\{u\leq-N\}).
\]
Choose $N$ with $a^n\mu(\{u\leq-N\})<V$. This is possible because $\mu$ is finite and gives no mass to the pole set. If $\sup_Xw_j<-aN$, its contact set lies in $\{u_j<-N\}$, contradicting the preceding inequalities and $\int_X\mathrm{MA}_\omega(w_j)=V$. Hence $\sup_Xw_j\geq-aN$ for $j\geq N$; for smaller $j$ the bound follows from $w_j\geq-aj$.

The envelopes decrease. Their uniform lower supremum prevents collapse to $-\infty$, by normalized quasi-plurisubharmonic compactness. Their genuine limit $w$ satisfies $w\leq au$, including at poles. With $\delta=1/a$, the function $\delta w$ is a genuine $\delta\omega$-plurisubharmonic competitor below $u$. Thus $P_{\delta\omega}(u)$ is genuine at every scale. The Lelong estimate preceding the theorem completes the proof. An additive normalization does not affect any conclusion.
\end{proof}

\begin{proposition}\label{prop:sublinear-threshold-envelope}\index{Envelope!sublinear singularity criterion}
Suppose a genuine quasi-plurisubharmonic $\psi\leq0$ satisfies $\ddc\psi\geq-C\omega$ for some $C>0$, and a genuine $\omega$-plurisubharmonic $u$ satisfies the canonical inequality
\[
 u\geq-A(1-\psi)^\alpha-B,
 \qquad A>0,\quad B\in\mathbb R,\quad0<\alpha<1.
\]
Then every $P_{\delta\omega}(u)$ is genuine, for $0<\delta<1$, and $u$ has zero Lelong numbers.
\end{proposition}
\begin{proof}
Choose $\alpha<\beta<1$ and $\varepsilon>0$ with $\varepsilon\beta C\leq\delta$. The number
\[
 K=B+\sup_{t\geq1}(At^\alpha-\varepsilon t^\beta)
\]
is finite, and $v_\delta=-\varepsilon(1-\psi)^\beta-K\leq u$. This composition is increasing and convex, with first derivative at most $\varepsilon\beta$ and nonnegative second derivative. Apply the quasi-plurisubharmonic chain inequality to $\max(\psi,-j)$ and pass to the decreasing limit. It gives $\ddc v_\delta\geq-\varepsilon\beta C\omega\geq-\delta\omega$. The limit is genuine because $\psi$ is finite almost everywhere. Its canonical pole values are $-\infty$, where the comparison still holds. Thus $v_\delta$ is an actual competitor at each scale. The preceding uniform Lelong estimate proves the last assertion.
\end{proof}

\subsubsection{A globally smooth root of rank at least two}

\begin{theorem}\label{thm:rank-two-dual-smooth-extension}\index{Gauduchon duality!smooth extension across degeneration}
Assume the semipositive attained setting of \cref{thm:semipositive-dual-exact-equation}, with greatest normalized threshold solution $\Phi_f$ and analytic degeneracy set $Z$. Suppose the maximizing form has the representation $T=\gamma^{n-1}$ for a globally smooth semipositive $(1,1)$-form $\gamma$ whose complex rank is at least two everywhere. If $\Phi_f$ has zero Lelong numbers, then it is smooth on $X$. This applies under either sufficient condition above. Every threshold subsolution is then smooth, since these subsolutions differ by constants.
\end{theorem}
\begin{proof}
Put $U=X\setminus Z$ and $A=\omega+\ddc\Phi_f$. The existing attained theorem makes $A$ smooth on $U$ and gives its exact global equation. The root $\gamma$ is positive on $U$. Stokes gives $\int_XA\wedge T=1$, while attainment gives
\[
 \int_U(A^n)^{1/n}(\gamma^n)^{(n-1)/n}
       =b_\omega(f)^{1/n}D_{f^{1/n}}(T)=1.
\]
The arithmetic--geometric mean gap is nonnegative and continuous on $U$, and its integral must vanish. Equality yields $A=r\gamma$ there for smooth $r\geq0$. Where $A^n=0$, equality forces $A=0$, so the same formula holds with $r=\operatorname{tr}_\gamma A/n=0$.

Since $\ddc\Phi_f$ is closed,
\[
 \mathrm dr\wedge\gamma+r\,\mathrm d\gamma=\mathrm d\omega
       \quad\hbox{on }U.
\]
For a semipositive form of rank at least two, the real-linear map $W_\gamma:\xi\mapsto\xi\wedge\gamma$ from one-forms to three-forms is injective. In a diagonal real unitary frame, for each coordinate of $\xi$ choose a positive coordinate two-plane not containing it. Its triple coefficient is that coordinate times the positive eigenvalue; no other disjoint diagonal pair contributes. Vanishing of the wedge therefore forces every coordinate to vanish.

Choose an auxiliary metric. The left inverse $L_\gamma=(W_\gamma^*W_\gamma)^{-1}W_\gamma^*$ is smooth and uniformly bounded on compact $X$. Thus
\begin{equation}\label{eq:rank-two-extension-first-order}
 \mathrm dr=\alpha-r\beta,\qquad
 \alpha=L_\gamma(\mathrm d\omega),\quad
 \beta=L_\gamma(\mathrm d\gamma),
\end{equation}
with smooth one-forms on $X$. In particular $|\mathrm d\log(1+r)|\leq C$ on $U$.

The analytic set permits short avoiding paths. In a coordinate ball, stratify it into countably many smooth strata of real dimension at most $2n-2$. For a fixed outside point, endpoints of segments meeting one stratum lie in a cone of dimension at most $2n-1$, a Lebesgue-null set. For two nearby outside points choose a generic intermediate point in a ball of radius comparable to their distance. Both connecting segments avoid $Z$ and their combined length is bounded by a fixed multiple of that distance. A finite connected coordinate cover, with outside points chosen in its overlaps, also gives a global uniform bound for connecting lengths in $U$.

The logarithmic gradient bound now bounds $r$ on $U$, and \eqref{eq:rank-two-extension-first-order} bounds its derivative. The local short paths give a unique locally Lipschitz extension across $Z$. Since $Z$ is Lebesgue-null, its weak derivative satisfies the same equation on $X$. The right side is continuous, so $r$ is $C^1$; iteration makes it smooth.

Set $S=r\gamma$. It is smooth semipositive and $\mathrm dS=\mathrm d\omega$ everywhere by density. The current $R=A-S$ is closed and supported on $Z$. It is positive: it equals the positive coefficient-measure part of $A$ on $Z$, since $A=S$ off $Z$ and $S$ has no mass there. The analytic support theorem recalled in \cref{sec:psh-divisorial-estimates} expresses $R$ as a nonnegative sum of divisorial integration currents; components of codimension at least two carry no such current. Locally $S-\omega$ is smooth and closed and has a smooth $\ddc$-potential. The divisorial coefficients are therefore the Lelong numbers of $\Phi_f$, and all vanish. Consequently $\ddc\Phi_f=S-\omega$ is smooth. Taking a Euclidean trace and applying elliptic regularity gives a smooth potential; its canonical representative agrees with it everywhere.
\end{proof}

The smooth root and its rank bound are explicit hypotheses. A smooth semipositive maximizing $(n-1,n-1)$-form need not have such a root across its degeneracy set. On a surface the rank bound already forces $\gamma>0$ everywhere.

\subsubsection{Two attained instances}

\begin{example}\label{ex:hopf-surface-invariant-attainment}\index{Gauduchon duality!Hopf surface attainment}
On the Hopf surface $X=(\mathbb C^2\setminus\{0\})/\langle z\mapsto2z\rangle$, the background $\omega=|z|^{-2}\ddc|z|^2$ has invariant bounded mass. It admits nonconstant smooth threshold solutions with actual smooth positive maximizing tests.
\end{example}
\begin{proof}
The form descends under $z\mapsto2z$ and is positive. For $s=|z|^2$ and $F(s)=s^{-1}$, the identity $sF''+2F'=0$ gives $\ddc\omega=\ddc(s^{-1})\wedge\ddc s=0$. It is nonclosed. Choose a nonconstant smooth real $q$ and small nonzero $\varepsilon$ such that $A=\omega+\ddc(\varepsilon q)>0$. Put $V=\int_X\omega^2$, $\mathrm{d}V=\omega^2/V$, $f=A^2/\mathrm{d}V$ and $u=\varepsilon q-\sup_X\varepsilon q$. Invariant mass gives $\int_Xf\,\mathrm{d}V=V$, hence $b_\omega(f)=1$ by \cref{thm:invariant-mass-threshold-envelopes}.

The test $T=A/V$ satisfies $\ddc T=0$, $\int_X\omega\wedge T=1$ and $D_{\sqrt f}(T)=1$. For every normalized smooth positive Gauduchon test $Q$, arithmetic--geometric mean and Stokes give $D_{\sqrt f}(Q)\leq\int_XA\wedge Q=1$. Thus $T$ attains the supremum. The attained uniqueness theorem identifies $\Phi_f=u$. This is a bounded nonconstant instance.
\end{proof}

\begin{example}\label{ex:quadric-rank-two-dual-attainment}\index{Gauduchon duality!quadric ramification}
For $n\geq3$, a smooth quadric double cover of $\mathbb P^n$ has an actual semipositive maximizing test whose smooth root has rank at least two everywhere and degenerates along the ramification divisor.
\end{example}
\begin{proof}
Let
\[
 X=\{[w:z_0:\cdots:z_n]:w^2=\textstyle\sum_{i=0}^nz_i^2\}
       \subset\mathbb P^{n+1},\qquad F([w:z])=[z].
\]
The quadratic is nonsingular and the map is finite of degree two. Its smooth ramification divisor is $Z=\{w=0\}$. Local branch coordinates show that $\mathrm dF$ has rank $n-1$ on $Z$ and $n$ elsewhere. Let $\omega$ be the restricted normalized Fubini--Study form and $A=F^*\omega_{\mathrm{FS}}$. Then
\[
 u=\log\frac{\sum_i|z_i|^2}{|w|^2+\sum_i|z_i|^2}
\]
is global and smooth, $-\log2\leq u\leq0$, $\sup_Xu=0$, and $A=\omega+\ddc u$. Branched change of variables gives $\int_XA^n=2$, and Stokes gives $V=\int_X\omega^n=2$.

Set $\mathrm{d}V=\omega^n/V$, $f=A^n/\mathrm{d}V$, $T=A^{n-1}/V$ and $\gamma=A/V^{1/(n-1)}$. The closed background gives $b_\omega(f)=1$. Stokes gives $\ddc T=0$ and $\int_X\omega\wedge T=1$; the determinant weights give $D_{f^{1/n}}(T)=1$. Arithmetic--geometric mean, valid also for semipositive $A$, bounds every competing normalized positive test by $\int_XA\wedge Q=1$. Thus $T$ is an actual maximizer and $\Phi_f=u$ by attained uniqueness. Its root has rank $n-1\geq2$ on $Z$ and is positive elsewhere, as required.
\end{proof}

The invariant-mass and sublinear conditions answer restricted instances of the envelope question. General moment or entropy bounds have not been shown to supply either condition or the smooth-root rank hypothesis.


\Needspace{14\baselineskip}
\begin{corollary}\label{cor:semipositive-dual-lp-limit}
Under the hypotheses of \cref{thm:semipositive-dual-exact-equation}, suppose in addition that $f\in L^p(\mathrm{d}V)$ for some $p>1$. Then the normalized threshold solution $u_0$ is continuous, and the normalized greatest exponential solutions converge uniformly to $u_0$ as the exponent tends to zero.
\end{corollary}
\begin{proof}
By \cref{prop:individual-normalizer-lp}, the constant $c$ for which $\mathrm{MA}_\omega(v)=cf\,\mathrm{d}V$ has a continuous solution is $b_\omega(f)$. The continuous normalized solution is therefore a threshold solution, and it equals $u_0$ by the uniqueness in \cref{thm:semipositive-dual-exact-equation}. For $L^p$ data the greatest exponential solution coincides with the bounded solution, as explained after \cref{thm:individual-exponential-subsolution}. By \cref{prop:symmetric-equation-clusters}, applied with the trivial group, the normalized family is uniformly precompact, and every cluster limit solves the threshold equation. Uniqueness identifies every cluster limit with $u_0$, which proves uniform convergence of the whole family.
\end{proof}


\subsection{A full-mass point pole without a smooth maximizing test}
\label{subsec:point-pole-dual-nonattainment}

The smaller-background criteria do not require $L^p$ integrability. The following closed-background example has genuine obstacle envelopes at every scale and an unbounded threshold solution. Its dual supremum is approached by smooth tests but is not attained by any smooth semipositive test.

\begin{proposition}\label{prop:point-pole-full-mass-nonattainment}\index{Gauduchon duality!point-pole nonattainment}\index{Monge--Amp\`ere density!point-pole integrability}
Let $X=\mathbb P^n$, $n\geq2$, and let $\omega$ be the normalized Fubini--Study form, with $\int_X\omega^n=1$ and $\mathrm{d}V=\omega^n$. For $0<\alpha<1$ put
\[
 \psi=\log\frac{\sum_{i=1}^n|z_i|^2}{\sum_{i=0}^n|z_i|^2},
 \qquad u=1-(1-\psi)^\alpha,
 \qquad o=[1:0:\cdots:0].
\]
The function $u$ is a genuine normalized $\omega$-plurisubharmonic potential with sole pole $o$ and extended mass one. Every $P_{\delta\omega}(u)$ is genuine, for $0<\delta<1$. Its measure is $f\,\mathrm{d}V$ with $f\in L^1(\mathrm{d}V)$ and $f\notin L^p(\mathrm{d}V)$ for every $p>1$.

One has $b_\omega(f)=1$ and $\rho_\omega(f^{1/n})=1$. Thus $u$ is an actual threshold solution. No smooth semipositive $T\in\mathscr G_\omega$ attains $D_{f^{1/n}}(T)=1$.
\end{proposition}
\begin{proof}
The numerator vanishes only at $o$. Locally $\omega+\ddc\psi=\ddc\log\sum_{i=1}^n|z_i|^2\geq0$. The function $\chi(t)=1-(1-t)^\alpha$ is increasing and convex on $t\leq0$, with $0\leq\chi'\leq\alpha$. Apply its chain inequality to bounded truncations of $\psi$ and take the decreasing limit. It gives a genuine $u\in\PSH(X,\omega)$, with
\[
 A=\omega+\ddc u\geq(1-\alpha)\omega.
\]
The function is smooth off $o$ and has supremum zero on $\{z_0=0\}$. Since $u=-(1-\psi)^\alpha+1$, \cref{prop:sublinear-threshold-envelope} gives all the smaller-background envelopes.

To compute its mass, use the chart $z_0=1$ and write $s=\sum_{i=1}^n|z_i|^2$. A local potential for $A$ and its radial flux are
\begin{equation}\label{eq:point-pole-radial-flux}
 \begin{aligned}
 F(s)&=\log(1+s)+1-\left(1-\log\frac{s}{1+s}\right)^\alpha,\\
 q(s)&=sF'(s)=\frac{s+\alpha(1-\log(s/(1+s)))^{\alpha-1}}{1+s}.
 \end{aligned}
\end{equation}
The tangential Hessian eigenvalue is $q/s$ and the radial one is $q'$. Polar integration gives the radial top-degree measure $nq^{n-1}q'\,\mathrm{d}s$. Its angular normalization follows from the Fubini--Study case $q=s/(1+s)$, whose radial density is $ns^{n-1}(1+s)^{-n-1}\,\mathrm{d}s$ and whose total mass is one. Consequently
\[
 \int_{\{a<|z|^2<b\}}A^n=q(b)^n-q(a)^n.
\]
Here $q(s)\to0$ as $s\downarrow0$ and $q(s)\to1$ as $s\to\infty$. On each compact annulus the extended measure equals the smooth top power. It gives no mass to $o$ or to the infinity hyperplane, since both are pluripolar. Exhausting the affine chart by annuli therefore gives extended mass exactly one. This is a mass calculation, independent of the vanishing of Lelong numbers.

Near $o$ write $L=1-\log s$. Direct differentiation of \eqref{eq:point-pole-radial-flux} gives
\[
 \begin{gathered}
 \frac qs\sim\alpha s^{-1}L^{\alpha-1},\qquad
 q'\sim\alpha(1-\alpha)s^{-1}L^{\alpha-2},\\
 f(s)\asymp s^{-n}L^{-1-n(1-\alpha)}.
 \end{gathered}
\]
The smooth radial volume is comparable to $s^{n-1}\,\mathrm{d}s$. Its $L^1$ tail is thus comparable to $\int^\infty L^{-1-n(1-\alpha)}\,\mathrm{d}L$, which is finite. For every $p>1$, the corresponding tail has an exponentially growing factor $\exp(n(p-1)L)$ and diverges. The measure is smooth off $o$, so it is globally $f\,\mathrm{d}V$. Since its total mass is one, \cref{thm:invariant-mass-threshold-envelopes} gives $b_\omega(f)=1$, and \cref{thm:integrable-gauduchon-formula} gives $\rho_\omega(f^{1/n})=1$.

Suppose a smooth semipositive normalized test $T$ attained this value. The positive current $A$ satisfies $\int_XA\wedge T=1$ by distributional Stokes. On $X\setminus\{o\}$, matrix arithmetic--geometric mean gives
\[
 A\wedge T\geq f^{1/n}k_T\,\mathrm{d}V.
\]
The right side has total mass one by attainment. Thus the nonnegative gap vanishes off $o$. There $A$ is smooth strictly positive, so equality gives $T=tA^{n-1}$, including at points where $T=0$. This equality case follows by applying arithmetic--geometric mean to the positive semidefinite coefficient matrix conjugated by $A^{1/2}$: all its eigenvalues agree, possibly at zero. The function $t=(T\wedge A)/A^n$ is smooth and nonnegative on the punctured manifold.

Since $A$ is closed, $\ddc T=0$ becomes $\ddc t\wedge A^{n-1}=0$, a locally strictly elliptic scalar equation there. All eigenvalues of $A$ tend to infinity at $o$, so, as $n\geq2$, all complementary-product eigenvalues of $A^{n-1}$ do also. Boundedness of the smooth coefficient matrix of $T$ forces $t\to0$ at $o$.

Extend $t(o)=0$. This continuous nonnegative function on compact $X$ is nonzero by the normalization of $T$. It has a positive maximum away from $o$. The local strong maximum principle on connected $X\setminus\{o\}$ forces it to be constant there, contrary to its zero limit at $o$. This excludes every smooth semipositive maximizing test; no assumption about an analytic degeneracy set was needed.
\end{proof}

\begin{corollary}\label{cor:point-pole-entropy-energy-threshold}\index{Entropy!point-pole threshold}
For the density and potential of \cref{prop:point-pole-full-mass-nonattainment},
\[
 \int_X f\log f\,\mathrm{d}V<+\infty
 \quad\Longleftrightarrow\quad \alpha<1-\frac1n.
\]
For every real $r>0$,
\[
 \int_X(-u)^r f\,\mathrm{d}V<+\infty
 \quad\Longleftrightarrow\quad \alpha r<n(1-\alpha).
\]
In particular, finite entropy of this density can coexist with an unbounded full-mass threshold solution and absence of a smooth semipositive maximizing test.
\end{corollary}
\begin{proof}
All possible divergence is near $o$. The preceding estimates give $f\,\mathrm{d}V\asymp L^{-1-n(1-\alpha)}\,\mathrm{d}L$, $\log f\sim nL$, and $-u\asymp L^\alpha$. The entropy tail converges exactly when $n(1-\alpha)>1$; the weighted tail converges exactly when $\alpha r<n(1-\alpha)$. At equality each tail is comparable to $\int^\infty L^{-1}\,\mathrm{d}L$ and diverges. The negative part of $f\log f$ is boundedly integrable on the probability volume, so the entropy criterion is unambiguous.
\end{proof}

These are entropy and weighted integrals of the single equation density. The original S16 question concerns uniform moment or entropy conditions for determinant-weighted probabilities of a maximizing sequence under smooth semipositive attainment. The present example lies outside that attained setting. The proposition concerns the explicit threshold solution $u$; its identification with the greatest threshold envelope would require a separate uniqueness argument.


\Needspace{19\baselineskip}
\subsection{Quantitative convergence of the currents}

The bound for the constants $c_\lambda$ in \cref{prop:semipositive-dual-exponential-limit} also bounds the trace error linearly in $\lambda$. Away from the degeneracy set, the deficit in the arithmetic--geometric mean inequality converts this bound into a rate of convergence for all coefficients of the currents. We keep the notation $u_0$, $v_\lambda$, $c_\lambda$ of that proposition. The total variation of a current of order zero is measured with respect to a fixed smooth Hermitian metric; on a relatively compact chart it is equivalent to the sum of the total variations of its coefficients.

\begin{proposition}\label{prop:semipositive-dual-current-rate}
Under the hypotheses and notation of \cref{prop:semipositive-dual-exponential-limit}, put
\[
 A_\lambda=\omega+\ddc v_\lambda,\qquad
 A_0=\omega+\ddc u_0.
\]
There is a constant $C$ such that
\[
 \|A_\lambda\wedge T-A_0\wedge T\|_{\mathrm{TV}}
       \leq C\lambda\qquad(\lambda>0).
\]
For each compact set $K\Subset X\setminus Z$, there is a constant $C_K$ such that
\[
 \|A_\lambda-A_0\|_{\mathrm{TV}(K)}
       \leq C_K\sqrt\lambda\qquad(0<\lambda\leq1).
\]
\end{proposition}
\begin{proof}
Put $p=f^{1/n}$, use the probability measure $\nu=A_0\wedge T=b^{1/n}p k_T\,\mathrm{d}V$ from the proof of \cref{prop:semipositive-dual-exponential-limit}, where the equality is \eqref{eq:semipositive-dual-trace-equality}, and set
\[
 q_\lambda=(c_\lambda/b)^{1/n}\mathrm{e}^{\lambda v_\lambda/n},
 \qquad R_\lambda=A_\lambda\wedge T-q_\lambda\nu\geq0.
\]
Here $R_\lambda\geq0$ by the trace inequality used in that proof, which also gives $c_\lambda\geq b$ and $\int_X(-v_\lambda)\,\mathrm{d}\nu\leq C_0$ for a constant $C_0$. Write $s_+=\max(s,0)$ for the positive part of a real number $s$. Since $q_\lambda\geq\mathrm{e}^{\lambda v_\lambda/n}$ and $1-\mathrm{e}^{x}\leq-x$, we obtain
\[
 (1-q_\lambda)_+\leq1-\mathrm{e}^{\lambda v_\lambda/n},
 \qquad
 \int_X(1-q_\lambda)_+\,\mathrm{d}\nu\leq C_0\lambda/n.
\]
Since $A_\lambda\wedge T$ and $\nu$ both have mass one,
\[
 \|A_\lambda\wedge T-\nu\|_{\mathrm{TV}}
 \leq R_\lambda(X)+\int_X|q_\lambda-1|\,\mathrm{d}\nu
 =2\int_X(1-q_\lambda)_+\,\mathrm{d}\nu.
\]
This proves the first estimate, with $C=2C_0/n$, and also gives
$R_\lambda(X)\leq C_0\lambda/n$ and
$\int_X|q_\lambda-1|\,\mathrm{d}\nu\leq2C_0\lambda/n$.

On $U=X\setminus Z$, write $T=\gamma^{n-1}$ and $A_0=a\gamma$ as in the proof of \cref{thm:semipositive-dual-exact-equation}, and put $\alpha_\lambda=q_\lambda a$. Decompose the positive current $A_\lambda$ into its absolutely continuous part, with matrix density $B_\lambda$ relative to $\gamma$, and its positive singular part. Since $\mathrm{MA}^{\mathrm{ext}}_\omega(v_\lambda)=c_\lambda\mathrm{e}^{\lambda v_\lambda}f\,\mathrm{d}V$, the trace inequalities for all Hermitian forms imply, as in the proof of \cref{cor:smooth-dual-exact-equation},
\[
 (\det B_\lambda)^{1/n}\geq\alpha_\lambda
 \quad\text{almost everywhere on }U.
\]
\Needspace{9\baselineskip}
This is a statement about determinants of matrices; no nonlinear product of coefficient measures is involved. Put $s_\lambda=\operatorname{tr}B_\lambda/n$ and $\delta_\lambda=s_\lambda-\alpha_\lambda\geq0$. On $U$, the absolutely continuous part of the trace error $R_\lambda$ is $\delta_\lambda\gamma^n$, and its singular part is the trace measure of the singular part of $A_\lambda$. Thus
\[
 \int_U s_\lambda\gamma^n\leq1,\qquad
 \int_U\delta_\lambda\gamma^n\leq C_0\lambda/n,
\]
and the singular trace mass on $U$ has the same upper bound.

For nonnegative numbers $\ell_1,\ldots,\ell_n$ with geometric mean at least $\alpha\geq0$ and arithmetic mean $s$, the arithmetic--geometric mean inequality and the Cauchy--Schwarz inequality give
\[
 \sum_i(\sqrt{\ell_i}-\sqrt\alpha)^2\leq n(s-\alpha),
 \qquad
 \sum_i|\ell_i-\alpha|\leq2n\sqrt{s(s-\alpha)}.
\]
Indeed $\sum_i\sqrt{\ell_i}\geq n\sqrt\alpha$, while
$\sum_i(\sqrt{\ell_i}+\sqrt\alpha)^2\leq4ns$.
Applying this to the eigenvalues of $B_\lambda$ relative to $\gamma$ and integrating, with the Cauchy--Schwarz inequality once more, we see that the absolutely continuous part of $A_\lambda-\alpha_\lambda\gamma$ has integrated trace norm at most $2n\sqrt{C_0\lambda/n}$ on $U$. Its singular part has norm at most $C_0\lambda$ in the same metric. Finally
\[
 \int_U|\alpha_\lambda-a|\gamma^n
   =\int_U|q_\lambda-1|\,\mathrm{d}\nu
   \leq2C_0\lambda/n.
\]
Since $\gamma$ is comparable to the fixed metric on the compact set $K\subset U$, these estimates and $\lambda\leq\sqrt\lambda$ for $0<\lambda\leq1$ give the second conclusion.
\end{proof}

The coefficient estimate is local on $X\setminus Z$. Convergence in capacity across $Z$ requires the additional envelope condition of \cref{cor:semipositive-dual-capacity-limit}.


\Needspace{16\baselineskip}
\subsection{The classes of the divisorial components}

Uniqueness constrains the divisors along which the threshold solution can have positive Lelong numbers. A linear relation between their Bott--Chern classes would produce a global potential that changes the coefficients of these divisors while preserving the equation off the degeneracy set, contradicting uniqueness. We write $\{D\}\in H^{1,1}_{\mathrm{BC}}(X,\mathbb R)$ for the Bott--Chern class of the current of integration $[D]$ along a prime divisor $D$, and $\nu(u,D)$ for the generic Lelong number of $u$ along $D$.

\begin{corollary}\label{cor:semipositive-dual-divisor-independence}
Under the hypotheses of \cref{thm:semipositive-dual-exact-equation}, let $u$ be a threshold solution and let $D_1,\ldots,D_k$ be the prime divisors with $\nu(u,D_j)>0$. Then the classes $\{D_j\}$ are linearly independent in $H^{1,1}_{\mathrm{BC}}(X,\mathbb R)$. In particular their number is at most $\dim_{\mathbb R}H^{1,1}_{\mathrm{BC}}(X,\mathbb R)$.
\end{corollary}
\begin{proof}
Since $u$ is smooth off $Z$, these divisors are among the finitely many components of $Z$ of codimension one. Put $a_j=\nu(u,D_j)>0$. By \cref{lem:nonclosed-divisorial-removal},
\[
 R=\omega+\ddc u-\sum_j a_j[D_j]\geq0.
\]
Suppose that real numbers $c_1,\ldots,c_k$ satisfy $\sum_j c_j\{D_j\}=0$. Fix a Gauduchon metric $\eta$ and apply \cref{lem:constant-trace-bc-representative} to the currents of integration $[D_j]$, with the integrable potentials $v_\eta([D_j])$ given there. Since the smooth representative in that lemma depends linearly on the Bott--Chern class, we obtain
\[
 h=\sum_j c_j v_\eta([D_j])\in L^1(X),\qquad
 \ddc h=\sum_j c_j[D_j].
\]
Choose $t\ne0$ sufficiently small that $a_j+tc_j\geq0$ for every $j$. The distribution $u_t=u+th$ satisfies
\[
 \omega+\ddc u_t=R+\sum_j(a_j+tc_j)[D_j]\geq0,
\]
so it is represented by a function $u_t\in\PSH(X,\omega)$. On $X\setminus Z$, the function $h$ is smooth and pluriharmonic, so $u_t$ has the same Monge--Amp\`ere measure as $u$ there. Since neither extended measure charges the pluripolar set $Z$, $u_t$ is another threshold solution. Uniqueness makes $th=u_t-u$ constant. Hence $\sum_jc_j[D_j]=0$, and testing near a generic point of each of the distinct divisors $D_j$ gives $c_j=0$.
\end{proof}


\Needspace{20\baselineskip}
\subsection{A branched projective example}

In this subsection $X=\mathbb P^n$, with integer $n\geq2$ and homogeneous coordinates $[Z_0:\cdots:Z_n]$. Let $\omega$ be the Fubini--Study form normalized by $\int_X\omega^n=1$, and put $\mathrm{d}V=\omega^n$. The current set $\mathscr G_\omega$, determinant density $k_T$, functionals $D_p$, $\rho_\omega$ and threshold $b_\omega$ have the meanings of \cref{sec:integrable-gauduchon-currents}.

A smooth semipositive maximizer need not be strictly positive, even when the threshold solution is smooth. The following example also shows why smoothness of the density root was asserted only off the degeneracy set.

\begin{example}\label{ex:semipositive-projective-maximizer}
For every $n\geq2$, there is a smooth nonnegative density $f$ on $\mathbb P^n$, vanishing precisely on the coordinate hyperplanes, such that $b_\omega(f)=1$ and $D_{f^{1/n}}$ has a smooth semipositive maximizer strictly positive off those hyperplanes. The threshold solution is smooth, but $f^{1/n}$ is not differentiable at a general point of any coordinate hyperplane. There is no smooth strictly positive maximizing form. The normalized greatest exponential solutions converge uniformly to the normalized threshold solution.
\end{example}
\begin{proof}
Define the holomorphic map $F:\mathbb P^n\to\mathbb P^n$ and the semipositive form $A$ by
\[
 F[Z_0:\cdots:Z_n]=[Z_0^2:\cdots:Z_n^2],\qquad
 A=\tfrac12F^*\omega.
\]
The homogeneous formula
\[
 u=\tfrac12\log\!\left(\sum_{i=0}^n|Z_i|^4\right)
       -\log\!\left(\sum_{i=0}^n|Z_i|^2\right)
\]
defines a smooth function on projective space, and $A=\omega+\ddc u$. This form is semipositive and is strictly positive exactly where every homogeneous coordinate is nonzero. Define
\[
 f=A^n/\omega^n,\qquad T=A^{n-1}.
\]
Closed Stokes gives $\int f\omega^n=\int\omega\wedge T=1$. Thus $T\in\mathscr G_\omega$. Every genuine $\omega$-plurisubharmonic potential has extended mass at most one: each bounded truncation has mass one, and restriction away from its truncation level can only decrease that mass. Consequently an inequality $\mathrm{MA}^{\mathrm{ext}}_\omega(v)\geq c f\omega^n$ forces $c\leq1$. The smooth solution $u$ gives the reverse bound, so $b_\omega(f)=1$.

The determinant formula for smooth semipositive forms gives
\[
 k_T=f^{(n-1)/n},\qquad
 D_{f^{1/n}}(T)=\int_{\mathbb P^n}f\omega^n=1.
\]
By \eqref{eq:integrable-current-supremum} and $\rho_\omega(f^{1/n})=b_\omega(f)^{-1/n}=1$, this is a maximizer. Its strict positivity away from the coordinate hyperplanes follows from that of $A$; on a hyperplane $A$ has deficient rank, and so does its cofactor matrix.

For the density root, use the chart $Z_0=1$ with coordinates $z_1,\ldots,z_n$. Pulling back the Fubini--Study volume by $(z_1,\ldots,z_n)\mapsto(z_1^2,\ldots,z_n^2)$ and including the factor $2^{-n}$ gives
\[
 f(z)=2^n\left(\prod_{i=1}^n|z_i|^2\right)
       \frac{(1+\sum_i|z_i|^2)^{n+1}}
            {(1+\sum_i|z_i|^4)^{n+1}}.
\]
Near a point of $\{z_1=0\}$ with all the other coordinates nonzero, its $n$th root is $|z_1|^{2/n}$ times a smooth strictly positive function. Restriction to a real line transverse to that hyperplane proves nondifferentiability, including the absolute-value case $n=2$. Coordinate symmetry gives the same conclusion on every component. The absence of a smooth strictly positive maximizer follows from \cref{cor:kahler-smooth-dual-attainment}, since $f$ has zeros. Finally $f$ is bounded, so \cref{cor:semipositive-dual-lp-limit} gives the uniform exponential limit.
\end{proof}


\subsection{A moment condition on maximizing tests}

For a general Hermitian background, a uniform bound on the moments of normalized potentials with respect to the probabilities of one smooth maximizing sequence can replace the regularity of the maximizing current. Together with the envelope condition in a smaller background, it shows that the cluster limits of the solutions of the exponential equations are exact threshold solutions.

For a nonzero $0\leq f\in L^1(\mathrm{d}V)$ with $b_\omega(f)>0$, let $\mathcal S_f$ be the class of threshold subsolutions with supremum zero and $\Phi_f$ its greatest member, from \cref{prop:greatest-normalized-threshold}. As in \cref{prop:semipositive-dual-exponential-limit}, $u_\lambda$ denotes the greatest solution of the exponential equation $\mathrm{MA}^{\mathrm{ext}}_\omega(u)=\mathrm{e}^{\lambda u}f\,\mathrm{d}V$, and $v_\lambda=u_\lambda-\sup_Xu_\lambda$, $c_\lambda=\mathrm{e}^{\lambda\sup_Xu_\lambda}$.

\Needspace{12\baselineskip}
\begin{theorem}\label{thm:gauduchon-moment-threshold}\index{Moment constraint}
Let $0\leq f\in L^1(\mathrm{d}V)$ be nonzero with $b=b_\omega(f)>0$, and put $p=f^{1/n}$ and $R=\rho_\omega(p)=b^{-1/n}$. Suppose that there are smooth strictly positive $T_j\in\mathscr G_\omega$, with associated probability measures $\nu_j$, such that
\[
 d_j=D_p(T_j)\longrightarrow R,\qquad
 \mathrm{d}\nu_j=d_j^{-1}p k_{T_j}\,\mathrm{d}V,
\]
and a constant $C$ satisfying
\begin{equation}\label{eq:maximizing-gauduchon-moments}
 \int_X(-v)\,\mathrm{d}\nu_j\leq C
 \quad\text{for all }j\text{ and all }v\in\PSH(X,\omega)
                  \text{ with }\sup_Xv=0.
\end{equation}
Then, for every $\lambda>0$,
\[
 b\leq c_\lambda\leq b\mathrm{e}^{C\lambda},\qquad
 \Phi_f-C\leq v_\lambda\leq0.
\]
Suppose in addition that $P_{\delta\omega}(\Phi_f)$ is not identically $-\infty$ for some $0<\delta<1$. Whenever $\lambda_k\downarrow0$ and $v_{\lambda_k}\to v$ in $L^1(X)$, the convergence holds in capacity and
\[
 \mathrm{MA}^{\mathrm{ext}}_\omega(v)=bf\,\mathrm{d}V,
 \qquad \Phi_f-C\leq v\leq\Phi_f.
\]
In particular an exact threshold solution exists.
\end{theorem}
\begin{proof}
Since $\mathrm{MA}^{\mathrm{ext}}_\omega(v_\lambda)=c_\lambda\mathrm{e}^{\lambda v_\lambda}f\,\mathrm{d}V$, the trace inequality of \cref{lem:unbounded-gauduchon-trace-density} and Stokes' theorem give, for each $j$ and $\lambda>0$,
\[
 1\geq c_\lambda^{1/n}d_j
              \int_X\mathrm{e}^{\lambda v_\lambda/n}\,\mathrm{d}\nu_j
   \geq c_\lambda^{1/n}d_j\mathrm{e}^{-C\lambda/n}.
\]
\Needspace{9\baselineskip}
The second inequality follows from Jensen's inequality and \eqref{eq:maximizing-gauduchon-moments}. Letting $j\to\infty$ and using $R=b^{-1/n}$, we obtain the upper bound for $c_\lambda$. The lower bound and the comparison with every normalized threshold subsolution were proved in \cref{cor:individual-positive-threshold-exponential}; applying that comparison to $\Phi_f$, we obtain
\[
 v_\lambda\geq\Phi_f-\lambda^{-1}\log(c_\lambda/b)
                \geq\Phi_f-C.
\]

Now let $v_{\lambda_k}\to v$ in $L^1$. By \cref{prop:greatest-normalized-threshold}, the limit $v$ lies in $\mathcal S_f$, hence below $\Phi_f$. The lower bound passes to the limit, so $\Phi_f-C\leq v\leq\Phi_f$; in particular $v\sim\Phi_f$ in the sense of \cref{def:singularity-order}. For all $k$ with $\lambda_k\leq1$,
\[
 \mathrm{MA}^{\mathrm{ext}}_\omega(v_{\lambda_k})
   =c_{\lambda_k}\mathrm{e}^{\lambda_k v_{\lambda_k}}f\,\mathrm{d}V
   \leq b\mathrm{e}^{C}f\,\mathrm{d}V.
\]
The functions $|\mathrm{e}^{v_{\lambda_k}}-\mathrm{e}^v|$ are uniformly bounded and tend to zero in $\mathrm{d}V$-measure, so their integrals against the fixed $L^1$ density $f$ tend to zero. Hence the error term in \eqref{eq:common-barrier-self-error} tends to zero, as required by \cref{prop:lower-barrier-criterion} with $\varphi=\Phi_f$. The additional hypothesis $v\preceq\Phi_f$ of that proposition holds since $v\leq\Phi_f$, and the condition on $P_{\delta\omega}(\Phi_f)$ is assumed. That proposition therefore gives convergence in capacity.

The continuity of damped measures under convergence in capacity, \cref{thm:damped-continuity}, now identifies the equation without a mass-conservation argument. Write $g_k=c_{\lambda_k}\mathrm{e}^{\lambda_k v_{\lambda_k}}f$. We know $g_k\to bf$ in $L^1$ and $0\leq g_k\leq b\mathrm{e}^Cf$ for large $k$. Hence
\[
 \mathrm{e}^{v_{\lambda_k}}\mathrm{MA}^{\mathrm{ext}}_\omega(v_{\lambda_k})
   =\mathrm{e}^{v_{\lambda_k}}g_k\,\mathrm{d}V
   \longrightarrow\mathrm{e}^v bf\,\mathrm{d}V
\]
in total variation, while \cref{thm:damped-continuity} gives the weak limit
$\mathrm{e}^v\mathrm{MA}^{\mathrm{ext}}_\omega(v)$. Dividing by $\mathrm{e}^v$ off the pole set of $v$, which neither measure charges, gives the equation and the finiteness of the mass. Finally, compactness of normalized $\omega$-plurisubharmonic functions provides an $L^1$ cluster limit along any sequence $\lambda_k\downarrow0$, which proves existence.
\end{proof}

The moment condition holds if the densities of the $\nu_j$ are uniformly bounded in $L^q(\mathrm{d}V)$ for some $q>1$, by H\"older's inequality and the uniform bounds for the moments of normalized $\omega$-plurisubharmonic functions. More generally, a uniform entropy bound suffices. Write $a_j=\mathrm{d}\nu_j/\mathrm{d}V$ and suppose $\int_Xa_j\log a_j\,\mathrm{d}V\leq H$ for all $j$. (By \cref{prop:hermitian-dual-entropy-bounds}, a uniform entropy bound is equivalent to a uniform bound for the dual entropy $H_0(\nu_j)$ with zero reference defined there.) Choose $\alpha>0$ and $A$ from \cref{lem:uniform-exponential}, so that $\int_X\mathrm{e}^{-\alpha v}\,\mathrm{d}V\leq A$ for every $v\in\PSH(X,\omega)$ with $\sup_Xv=0$. The elementary entropy inequality, applied first to the truncations $v_T=\max(v,-T)$ at levels $T>0$, gives
\[
 \alpha\int_X(-v_T)a_j\,\mathrm{d}V
 \leq \int_Xa_j\log a_j\,\mathrm{d}V
             +\log\int_X\mathrm{e}^{-\alpha v_T}\,\mathrm{d}V
 \leq H+\log A.
\]
To prove this inequality, apply $t\log t-t+1\geq0$ to $t=a_j/q$, integrate against $q\,\mathrm{d}V$, where $q=\mathrm{e}^{-\alpha v_T}/\int_X\mathrm{e}^{-\alpha v_T}\,\mathrm{d}V$ is a probability density, and rearrange. Monotone convergence as $T\to\infty$ proves \eqref{eq:maximizing-gauduchon-moments} with $C=(H+\log A)/\alpha$.

The same entropy bound also gives attainment of the dual supremum along this sequence. Indeed, $t\log t\geq-\mathrm{e}^{-1}$ and, for $K>1$,
\[
 \int_{\{a_j>K\}}a_j\,\mathrm{d}V
       \leq\frac{H+\mathrm{e}^{-1}}{\log K}.
\]
This proves uniform integrability of the $a_j$, and since $d_j\to R>0$, also of the densities $p k_{T_j}=d_ja_j$. By \cref{prop:selected-gauduchon-attainment}, every weak cluster limit $T$ of $(T_j)$ maximizes $D_p$, and along a subsequence converging to $T$ the probabilities $\nu_j$ converge to $R^{-1}p k_T\,\mathrm{d}V$. For each normalized $v$, the function $-v$ is nonnegative and lower semicontinuous, so the lower semicontinuity of its integral under weak convergence of measures passes the moment bound to this limiting probability. These conclusions give no regularity of $T$.

The envelope condition holds if $\Phi_f$ is bounded. In general, it and the moment bound are additional hypotheses; the tail condition on $f^{1/n}$ alone does not provide them in this argument. The theorem identifies the equation satisfied by each cluster limit. Uniqueness, which gave convergence of the whole family in \cref{prop:semipositive-dual-exponential-limit}, is not available here.


